\documentclass[10pt]{article}

\usepackage[margin=0.75in]{geometry}
\usepackage{amsmath}
\usepackage{booktabs}
\usepackage{caption}
\usepackage{graphicx}
\usepackage{hyperref}
\usepackage{longtable}
\usepackage{makecell}
\usepackage{pdflscape}
\usepackage{tabularx}

\hypersetup{colorlinks=true, linkcolor=blue, urlcolor=blue}
\setlength{\parindent}{0pt}
\setlength{\parskip}{0.6em}
\renewcommand{\arraystretch}{1.12}

\newcommand{\sym}[1]{#1}
\newcommand{\todosection}[2]{\clearpage\section*{#1. #2}\addcontentsline{toc}{section}{#1. #2}}
\newcommand{\feedback}[1]{\textbf{Feedback.} #1\par}
\newcommand{\response}[1]{\textbf{What we added or changed.} #1\par}
\newcommand{\readthis}[1]{\textbf{How to read the exhibit.} #1\par}
\newcommand{\inputtable}[1]{\IfFileExists{#1}{\resizebox{\linewidth}{!}{\input{#1}}}{\emph{Table file not found.}}}

\begin{document}

\begin{center}
{\Large ToDo Reviews}\\
\vspace{0.35em}
Summary of feedback, paper responses, and changed tables/figures
\end{center}

\tableofcontents

\section*{Overview}
\addcontentsline{toc}{section}{Overview}

Each section summarizes one referee-feedback item, explains what we added or changed in the paper, and shows the table or figure affected by that change when there is one. Items that only required text changes are summarized without code or commit details.

\todosection{1}{Prior Deworming Balance}
\feedback{The referee worried that baseline imbalance in prior deworming could explain the stronger public-signal effects in Far communities.}
\response{We added an appendix robustness table that controls for cluster-level baseline shares of adults who had ever dewormed and who had dewormed in the past year. The paper text also explains why the Calendar arm is a useful placebo: Calendar has similar prior-deworming imbalance patterns but does not show the same Far--Close take-up pattern as the public signals.}
\readthis{The key comparison is the Far--Close column. Bracelet and Ink remain positive after controlling for prior-deworming shares, while Calendar remains small.}

\begin{landscape}
\begin{table}[p]
\centering
\caption*{Prior-deworming robustness}
\inputtable{presentations/rf-tables/main-specs/prior-deworming-robustness.tex}
\end{table}
\end{landscape}

\todosection{2}{Balance Table Control Means}
\feedback{The referee noted that some balance-table ``Control mean'' values looked mechanically inconsistent with the overall sample means and treatment differences.}
\response{We fixed the balance tables so the displayed Control means are raw Control-group means in the relevant sample, not fixed-effect regression intercepts for the omitted county stratum.}
\readthis{The corrected table now reports intuitive Control means. For example, in the take-up sample, Age is about 40.3 in Control-Close and 40.2 in Control-Far; Female is 0.565 and 0.562; Phone owner is 0.765 and 0.728.}

\begin{landscape}
\begin{table}[p]
\centering
\caption*{Corrected balance table}
\inputtable{presentations/tables/stacked-sample-balance-table.tex}
\end{table}
\end{landscape}

\begin{landscape}
\begin{table}[p]
\centering
\caption*{Corrected distance-balance table}
\inputtable{presentations/tables/stacked-sample-dist-balance-table.tex}
\end{table}
\end{landscape}

\todosection{4}{Units for the Relative Valuation Parameter}
\feedback{The referee pointed out that the valuation parameter \(\psi\) appeared to mix KSh and USD units.}
\response{We changed the paper and generated tables so \(\psi\), the mean calendar-minus-bracelet valuation difference, is reported on the USD scale.}
\readthis{The exhibit shows only the affected row. The estimate did not change; the substantive change is the unit label for the valuation difference.}

\begin{table}[h!]
\centering
\caption*{Relative valuation row before and after the unit-label correction}
\begin{tabular}{lll}
\toprule
Version & Parameter row & Estimate\\
\midrule
Old & \(\psi\) & \makecell[l]{0.472\\(0.403, 0.547)}\\
New & \(\psi\text{ (USD)}\) & \makecell[l]{0.472\\(0.403, 0.547)}\\
\bottomrule
\end{tabular}
\end{table}

\todosection{5 and 8}{Continuous Distance and Assignment Simulations}
\feedback{The referee wanted a clearer explanation of why realized continuous distance can be interpreted as design-induced variation, given the constrained PoT selection and assignment process.}
\response{We expanded the paper to explain that the simulation checks hold the surveyed-community frame fixed and rerun the constrained PoT selection, Close/Far assignment, targeted-community pairing, and feasibility screening.}
\readthis{The figure shown is the new design-matched counterfactual assignment density used in the paper. It uses counterfactual distance assignments generated by rerunning the experiment's constrained assignment procedure.}

\begin{figure}[p]
\centering
\includegraphics[width=0.72\linewidth]{presentations/refine-review-artifacts/new-pot-distance-overlap.pdf}
\caption*{New potential-distance plot using counterfactual distance assignments from the experiment assignment procedure}
\end{figure}

\todosection{6}{Identification of \(\sigma_u\)}
\feedback{The referee wanted a clearer account of what disciplines the idiosyncratic-noise parameter \(\sigma_u\), and whether it is mainly coming from the prior.}
\response{We added text explaining that \(\sigma_u\) is informed by the curvature of take-up across arm-by-distance cells and by the normal-sum restriction linking intrinsic motivation, decision noise, and observed deworming choices. We also added an appendix figure comparing the prior and posterior distributions for \(\sigma_u\).}
\readthis{The posterior is shifted toward smaller values than the prior, which is consistent with the data implying that deworming choices are informative about intrinsic motivation.}

\begin{figure}[p]
\centering
\includegraphics[width=0.72\linewidth]{presentations/figures/sigma-u-prior-posterior.pdf}
\caption*{Prior and posterior distributions for \(\sigma_u\)}
\end{figure}

\begin{table}[h!]
\centering
\caption*{Summary of prior and posterior for \(\sigma_u\)}
\begin{tabular}{lrrrr}
\toprule
Distribution & Mean & Median & 2.5\% & 97.5\%\\
\midrule
Prior & 0.468 & 0.374 & 0.145 & 1.359\\
Posterior & 0.302 & 0.278 & 0.127 & 0.583\\
\bottomrule
\end{tabular}
\end{table}

\todosection{10}{Figure 3(b) Curve Ordering}
\feedback{The referee suggested that Figure 3(b) might show an impossible crossing of curves.}
\response{We checked the plotted data and confirmed that the highlighted curves do not cross on their common support. We therefore did not add a final paper change for this item.}
\readthis{No table or figure changed.}

\todosection{11}{HMC Diagnostics}
\feedback{The referee thought the original Bayesian convergence statement, split \(\hat R<1.1\), was too loose for the structural model.}
\response{We replaced the broad statement with model-specific diagnostics in the paper: maximum split \(\hat R\), minimum effective sample sizes, and divergent-transition counts for the main structural fit and appendix robustness fits.}
\readthis{The main model and most robustness fits have no divergences. One full-information robustness fit has two divergent transitions, which is now reported explicitly.}

\begin{longtable}{p{0.24\linewidth}p{0.58\linewidth}p{0.12\linewidth}}
\toprule
Fit/table family & Diagnostics reported & Status\\
\midrule
Main structural fit and downstream WTP/ATE/optimization tables & Max split \(\hat R=1.011\), min bulk ESS 465, min tail ESS 602, zero divergent transitions. & Added\\
No-outlier structural appendix & Max split \(\hat R=1.018\), min bulk ESS 355, min tail ESS 309, zero divergent transitions. & Added\\
Community-FP individual-distance appendix & Max split \(\hat R=1.022\), min bulk ESS 221, min tail ESS 212, zero divergent transitions. & Added\\
Full-information individual-distance appendix & Max split \(\hat R=1.022\), min bulk ESS 142, min tail ESS 132, two divergent transitions. & Added\\
\bottomrule
\end{longtable}

\todosection{12}{Continuous-Distance Balance Test}
\feedback{The referee noted that the old permutation test for continuous distance did not obviously match the constrained assignment design.}
\response{We regenerated the E3 balance test using feasible counterfactual distance draws from the assignment simulations rather than freely permuting observed distances within counties.}
\readthis{The new plot is the design-matched version. The conclusion remains that realized distance is not systematically related to the displayed covariates.}

\begin{figure}[p]
\centering
\begin{minipage}{0.48\linewidth}
\centering
\includegraphics[width=\linewidth]{presentations/refine-review-artifacts/old-dist-ri-plot.pdf}
\caption*{Old E3 plot}
\end{minipage}
\hfill
\begin{minipage}{0.48\linewidth}
\centering
\includegraphics[width=\linewidth]{presentations/refine-review-artifacts/new-dist-ri-plot.pdf}
\caption*{New E3 plot}
\end{minipage}
\end{figure}

\todosection{17}{Figure M1 Parameters}
\feedback{The referee noted that the text describing Figure M1 did not match the distributions actually plotted.}
\response{We updated the appendix text and figure note so they match the plotted distributions: a narrow Gaussian benchmark around 0.4 and an asymmetric bimodal distribution with one mode near -2 and one near 0.4.}
\readthis{This was a text/caption correction; no table changed.}

\todosection{18}{Closest PoT and Original-Distance ITT}
\feedback{The referee saw a tension between the design statement about being closest to the assigned PoT and the implementation discussion of finalized PoT locations.}
\response{We clarified the distinction between candidate PoT centers used in the assignment algorithm and finalized implemented PoT locations. We also verified in the realized implementation data that each selected community's assigned finalized PoT is the closest finalized deworming site. As an additional robustness check, we added original-distance-assignment ITT tables for take-up and first-order observability.}
\readthis{The ITT tables use the original randomized Close/Far distance assignment. The take-up and observability patterns remain broadly consistent with the main results.}

\begin{landscape}
\begin{table}[p]
\centering
\caption*{Original-distance-assignment ITT: take-up}
\inputtable{presentations/rf-tables/main-specs/rf_original_distance_itt_tbl.tex}
\end{table}
\end{landscape}

\begin{landscape}
\begin{table}[p]
\centering
\caption*{Original-distance-assignment ITT: first-order observability}
\inputtable{presentations/rf-tables/main-specs/rf_original_distance_itt_fob_tbl.tex}
\end{table}
\end{landscape}

\todosection{24}{Normal-Sum Terminology}
\feedback{The referee noted that ``normal mixture'' was statistically imprecise because the model uses a sum of normal variables, not a mixture distribution.}
\response{We changed the wording to ``normal-sum specification.''}
\readthis{Text-only correction.}

\begin{quote}
Old: ``For the normal mixture, the type inference term is ...''

New: ``Under this normal-sum specification, the type inference term is ...''
\end{quote}

\todosection{25}{SMS Control-Mean Standard Error and Table Style}
\feedback{The referee flagged that the SMS table appeared to report a standard deviation below the control mean, while the paper intended to report a standard error.}
\response{We changed the table generation so the SMS control mean parenthetical reports a clustered standard error. We also made the SMS and heterogeneity table output reproducible in a style closer to the manually edited paper tables.}
\readthis{The SMS control mean is 0.330 and the parenthetical is now 0.032, replacing the old 0.470 standard deviation.}

\begin{table}[h!]
\centering
\caption*{SMS treatment effects after standard-error fix}

\begingroup
\centering
\begin{tabular}{lc}
   \tabularnewline \midrule \midrule
   Model:        & (1)\\  
   \midrule
   \emph{Variables}\\
   Reminder Only & 0.167$^{***}$\\   
                 & (0.029)\\   
   Social Info   & 0.131$^{***}$\\   
                 & (0.028)\\   
   \midrule
   \emph{Fit statistics}\\
   Control mean  & 0.330\\  
                 & (0.032)\\  
   Observations  & 1,705\\  
   \midrule \midrule
\end{tabular}
\par\endgroup



\end{table}

\begin{landscape}
\begin{table}[p]
\centering
\caption*{Regenerated heterogeneity table}
\inputtable{presentations/rf-tables/main-specs/het-tes.tex}
\end{table}
\end{landscape}

\end{document}
